\subsection*{2.4 Borel Sets}

In this section, we demonstrate how to construct a $\sigma$-field that contains a given collection of subsets in $\Omega$. In particular, suppose we are interested in a topological space. The open sets in the topological space may not satisfy all the requirements of a $\sigma$-algebra. In this case, we define the \textit{Borel algebra} as the smallest $\sigma$-field that contains all open sets in the given topological space.

\begin{thmbox}{2.6}
Let $\mathcal{C}$ be a collection of subsets in $\Omega$, and $\mathcal{F}$ be the intersection of all $\sigma$-fields that contain $\mathcal{C}$ as a sub-collection. Then $\mathcal{F}$ is a $\sigma$-field, and it is the smallest $\sigma$-field that contains $\mathcal{C}$.
\end{thmbox}

\textit{Proof} We can write $\mathcal{F}$ as
\[
\mathcal{F}
=
\bigcap_{\mathcal{H}:\sigma\text{-field},\ \mathcal{C}\subseteq \mathcal{H}} \mathcal{H},
\]
where the intersection taken over all $\sigma$-fields that contain $\mathcal{C}$ as a sub-collection. The intersection is not empty because $\mathcal{P}(\Omega)$ is such a $\sigma$-field. We next check that $\mathcal{F}$ is a $\sigma$-field.

To show that $\mathcal{F}$ is a $\sigma$-field, we need to verify the three axioms:
\begin{itemize}
    \item $\Omega\in \mathcal{F}$ because $\Omega$ is an element of every $\sigma$-field, in particular, every $\sigma$-field that contains $\mathcal{C}$.
    \item If $A\in \mathcal{F}$, then $A$ is contained in every $\sigma$-field that contains $\mathcal{C}$. Since any $\sigma$-field is closed under taking complements, $A^c$ is also contained in every $\sigma$-field that contains $\mathcal{C}$. This proves that $A^c\in \mathcal{F}$.
    \item Suppose $A_i\in \mathcal{F}$, for $i=1,2,\dots$. By definition, $A_i$ is contained in every $\sigma$-field that contains $\mathcal{C}$, so their union $\cup_i A_i$ is also contained in every $\sigma$-field containing $\mathcal{C}$. Therefore $\cup_i A_i$ is an element of $\mathcal{F}$.
\end{itemize}

This proves that $\mathcal{F}$ is a $\sigma$-field that contains $\mathcal{C}$. Moreover, if $\mathcal{G}$ is any $\sigma$-field that contains $\mathcal{C}$, then $\mathcal{G}$ is one of the $\sigma$-fields being intersected to form $\mathcal{F}$. Therefore, $\mathcal{F}$ is a subset of $\mathcal{G}$. This implies that $\mathcal{F}$ is the smallest $\sigma$-field containing $\mathcal{C}$. \hfill $\square$

\begin{defbox}{2.7}
The $\sigma$-field $\mathcal{F}$ in Theorem 2.6 is called the $\sigma$-\textit{field generated by} $\mathcal{C}$ and is denoted by $\sigma(\mathcal{C})$. If $(B_i)_{i\ge 1}$ is a sequence of sets in $\Omega$, then the $\sigma$-field generated by this sequence of sets is denoted by $\sigma((B_i)_{i\ge 1})$.
\end{defbox}

\textbf{Example 2.4.1 (Generating a $\sigma$-algebra in a Finite Measure Space)} \\
Consider a sample space $\Omega=\{a,b,c,d,e,f\}$. Let $A_1=\{a,b,c\}$ and $A_2=\{c,d,e\}$ be two of its subsets. We aim to list all the sets in the $\sigma$-algebra generated by $\{A_1,A_2\}$. Since the sample space is finite, a $\sigma$-algebra is the same as an algebra.

Let $\mathcal{F}$ denote the algebra generated by $A_1$ and $A_2$. We will list the subsets of $\mathcal{F}$ in pairs, along with their complement. First, note that $\mathcal{F}$ must contain the empty set, $A_1$, $A_2$, and their complements. Therefore, the following sets
\[
\emptyset,\ \{a,b,c,d,e,f\},
\]
\[
A_1=\{a,b,c\}, \quad A_1^c=\{d,e,f\},
\]
\[
A_2=\{c,d,e\}, \quad A_2^c=\{a,b,f\},
\]
are in $\mathcal{F}$.

Because $\mathcal{F}$ is an algebra and is closed under taking union and intersection, it must contain the union and intersection of any pair of the subsets above and their complements. Therefore, the following sets should also belong to $\mathcal{F}$:
\[
\{a,b,c,d,e\},\ \{f\}, \qquad \{c\},\ \{a,b,d,e,f\},
\]
\[
\{a,b,c,f\},\ \{d,e\}, \qquad \{a,b\},\ \{c,d,e,f\}.
\]

We have listed $14$ subsets so far, and we expect that there the number of subsets in $\mathcal{F}$ is a power of $2$. Indeed, the last two subsets in $\mathcal{F}$ are
\[
\{a,b,d,e\},\ \{c,f\}.
\]

Note that $\{a,b,d,e\}$ is the symmetric difference of $A_1$ and $A_2$. One can check that these subsets are closed under union, intersection, and complement and hence form an algebra. This is the smallest algebra that contains $A_1$ and $A_2$.

\begin{defbox}{2.8}
Let $\Omega=\mathbb{R}$ and $\mathcal{C}$ be the collection of all open intervals $(a,b)$, with $a<b$. The $\sigma$-field generated by $\mathcal{C}$ is a $\sigma$-field called the \textit{Borel field}, or the \textit{Borel algebra}, and is written as $\mathcal{B}(\mathbb{R})$. Likewise, for integer $d\ge 1$, we define the Borel algebra $\mathcal{B}(\mathbb{R}^d)$ as the $\sigma$-algebra generated by the open balls in $\mathbb{R}^d$. A set in a Borel algebra is called a \textit{Borel set}.

The Borel algebra $\mathcal{B}(\mathbb{R})$ contains all subsets that we need in practice. For example, closed intervals $[a,b]$ are in $\mathcal{B}(\mathbb{R})$ because it can be expressed as an intersection
\[
[a,b]=\bigcap_{n=1}^{\infty} \left(a-\frac{1}{n},\, b+\frac{1}{n}\right).
\]

Semi-infinite intervals $(-\infty,b]$ are in $\mathcal{B}(\mathbb{R})$ because
\[
(-\infty,b]=\bigcup_{n=1}^{\infty} [b-n,b].
\]

Similarly, $[a,\infty)$, $(a,\infty)$, $(-\infty,b)$, $(a,b]$, $[a,b)$ are Borel sets.
\end{defbox}

In Definition 2.8, we generate the Borel sets by open intervals. Nevertheless, there is nothing special about open intervals, and we can use other types of intervals instead.

\begin{thmbox}{2.7}
The Borel algebra on $\mathbb{R}$ can be generated by closed intervals $[a,b]$, for $a<b$.
\end{thmbox}

\textit{Proof} Let $\mathcal{C}$ denote the collection of open intervals, and let $\mathcal{C}'$ be the collection of closed intervals in $\mathbb{R}$. We want to show that $\sigma(\mathcal{C})=\sigma(\mathcal{C}')$. This is equivalent to prove (i) $\sigma(\mathcal{C})\subseteq \sigma(\mathcal{C}')$ and (ii) $\sigma(\mathcal{C})\supseteq \sigma(\mathcal{C}')$.

By the definition of $\sigma(\mathcal{C}')$, we have
\[
\sigma(\mathcal{C}')
=
\bigcap_{\mathcal{H}:\sigma\text{-field},\ \mathcal{C}'\subseteq \mathcal{H}} \mathcal{H}.
\]

As we have shown in the paragraph before the theorem, the $\sigma$-algebra $\sigma(\mathcal{C})$ contains all closed intervals and is one of the $\sigma$-algebras $\mathcal{H}$ in the intersection. This proves that $\sigma(\mathcal{C}')\subseteq \sigma(\mathcal{C})$.

In the other direction, the $\sigma$-algebra $\sigma(\mathcal{C}')$ contains all open intervals because an open interval $(a,b)$ can be written as
\[
(a,b)=\bigcup_{n=1}^{\infty} \left[a+\frac{1}{n},\, b-\frac{1}{n}\right].
\]
Therefore $\sigma(\mathcal{C})\subseteq \sigma(\mathcal{C}')$. \hfill $\square$

Similarly, one can show that $\mathcal{B}(\mathbb{R})$ can be generated by sets in the form $(a,b]$, or sets in the form $(-\infty,b]$, etc.

As in the one-dimensional case, $\mathcal{B}(\mathbb{R}^n)$ can be generated by other classes of sets, such as closed $n$-dimensional boxes, or open $n$-dimensional balls.

\begin{thmbox}{2.8}
For positive integer $n$, the Borel algebra on $\mathbb{R}^d$ is generated by open $d$-dimensional open boxes in the form
\[
(a_1,b_1)\times (a_2,b_2)\times \cdots \times (a_d,b_d).
\]
\end{thmbox}

The proof is similar to the proof of Theorem 2.7. We just need to show that an open disc in $\mathbb{R}^d$ can be expressed as a countable union of open $d$-dimensional boxes and \textit{vice versa}.
